\documentclass[11pt,a4paper]{article}

\usepackage[margin=2.5cm]{geometry}
\usepackage{amssymb}
\usepackage{booktabs}
\usepackage{tabularx}
\usepackage{multirow}
\usepackage{hyperref}
\usepackage[table]{xcolor}
\usepackage{enumitem}
\usepackage{lmodern}
\usepackage[T1]{fontenc}
\usepackage{microtype}

\hypersetup{
  colorlinks=true,
  linkcolor=blue!60!black,
  urlcolor=blue!60!black,
}

\setlength{\parindent}{0pt}
\setlength{\parskip}{0.6em}

% Long \texttt{} identifiers (mnemonic names, parameter set strings)
% are unbreakable single words. Give TeX extra stretch budget to find
% acceptable line breaks at the surrounding spaces.
\sloppy
\emergencystretch=3em

\title{Re-implementing XMSS for Verification:\\
  A Project Summary}
\author{Wrenna Robson\thanks{With Claude (Anthropic), which co-authored
  the implementations, analysis tooling, and report text. The work
  presented here is intended for evaluation and review and is
  \emph{not} approved for production use.}}
\date{April 2026}

\begin{document}
\maketitle

\begin{abstract}
This report summarises a multi-implementation re-engineering of the
XMSS and XMSS-MT stateful hash-based signature schemes (RFC~8391,
including Errata~7900). Two implementations were produced from
scratch: a C99 reference deliberately written to a restrictive
subset that mirrors the constraints of the Jasmin language, and a
Jasmin implementation covering five SHA-2/256 parameter sets (three
single-tree and two multi-tree) with constant-time guarantees
verified by the \texttt{jasmin-ct} checker on x86-64. The report
profiles the RISC-V ISA surface of both implementations, develops a
scoping analysis for a future Jasmin RISC-V port, and presents a
candid account of currently known limitations. RISC-V ISA-specific
optimisation work for XMSS is concentrated almost entirely in the
hash layer: the algorithm layer compiles to portable RV64I + M and
benefits negligibly from any extension beyond it.
\end{abstract}

\section{Context and goals}

XMSS is a stateful hash-based signature scheme standardised in
RFC~8391 and approved by NIST in SP~800-208 as a post-quantum
algorithm for stateful use cases such as firmware and code signing.
Its security reduces to assumptions on the underlying hash function
alone, and, unlike lattice schemes, it does not depend on
algebraic structure that may be amenable to future quantum or
classical cryptanalysis. The trade-off is statefulness: each secret
key may sign only a fixed number of messages, and the signing state
must be persisted and incremented atomically across signing
operations.

This project pursues two goals that interact:

\begin{enumerate}[nosep]
  \item Produce an XMSS implementation suitable for formal
    verification of constant-time properties and (eventually)
    functional correctness, in a language whose semantics support
    such proofs. The chosen target is
    \href{https://github.com/jasmin-lang/jasmin}{Jasmin}, a
    restricted assembly-like language whose compiler preserves a
    constant-time discipline checkable at the source level by the
    \texttt{jasmin-ct} tool, and whose semantics are formalised in
    EasyCrypt.
  \item Inform a future RISC-V~64 port. The Jasmin x86-64 backend
    is mature; the RISC-V backend is under active development. We
    want to know, before committing to that port, what ISA surface
    XMSS exercises and where work on the backend (or on
    architecture-specific code generation) would pay off.
\end{enumerate}

The C~implementation is not a competitor to the upstream reference
maintained by the RFC authors; its purpose is to serve as a
Jasmin-shaped specification. Every choice that diverges from the
reference is motivated by what the corresponding Jasmin code must
look like, given Jasmin's prohibition of variable-length arrays,
heap allocation, function pointers, and recursion. The result is a
C codebase whose structure---module boundaries, calling conventions,
buffer sizes, address-serialisation discipline---is a direct preview
of the Jasmin code, and against which the Jasmin code can be
cross-validated byte for byte.

\section{Scope of the implementations}

\subsection{C99 reference}

The C implementation supports all twelve XMSS and thirty-two XMSS-MT
parameter sets registered in RFC~8391, across SHA-2/256, SHA-2/512,
SHAKE-128, and SHAKE-256 hash backends. It uses the BDS authentication
path algorithm \cite{BDS09} for amortised $O(1)$ signing cost.

The library is allocation-free: every buffer is either stack-local
or caller-provided. The implementation respects an enforced subset
that prohibits VLAs, function pointers, heap allocation, recursion,
and any secret-dependent control flow or memory access. These
constraints are tighter than RFC~8391 requires, and tighter than the
upstream reference observes; they exist solely to make the code
structurally portable to Jasmin.

Test coverage spans unit tests of every module, sign/verify
roundtrips for every supported parameter set, BDS state
serialisation, XMSS-MT tree boundary crossings, KAT cross-validation
against the upstream reference, ACVP cross-checks, and exhaustive
sign/verify across every leaf index for valid small-height
parameter combinations.

\subsection{Jasmin implementation}

The Jasmin implementation covers five parameter sets at SHA-2/256:
three single-tree and two multi-tree. Each parameter set is exposed
as its own compiled assembly module with libjade-style export
naming. Keygen, sign, and verify are tested for every set; the
multi-tree sets are tested through tree boundary crossings (more
than 1024 signatures).

The implementation passes the \texttt{jasmin-ct} constant-time
checker on x86-64 across all twenty-four exported entry points. A
small number of declassification sites are present and individually
justified; the principal one is the leaf index, which is secret
until written into the signature and public thereafter.

Cross-implementation interop tests sign with both the C and Jasmin
backends from the same seed and compare both signatures and BDS
state byte-for-byte after every signature, accounting for the
endianness difference between the C wire-format big-endian state
words and the Jasmin native little-endian internal representation.

\section{Re-implementing XMSS in C: structural choices}

The C implementation makes several design choices that visibly
inflate its instruction count relative to the upstream reference.
These are deliberate preparation for the Jasmin port. The following
module-level deltas (measured on RV64 with \texttt{-O3}; full data
in the comparative ISA report) are representative:

\begin{table}[ht]
\centering
\caption{Selected per-module instruction count differences between
  this project's C implementation and the upstream reference, with
  the design reason for each.}
\label{tab:tax}
\begin{tabularx}{\linewidth}{@{}lrX@{}}
\toprule
Module                & Delta & Source of the difference \\
\midrule
Parameter derivation  & $-$791 & Compact OID lookup table replacing a switch cascade (a Jasmin-neutral C-level optimisation). \\
Hash wrappers         & $+$398 & ADRS pre-serialised to a stack buffer and passed by pointer, matching the Jasmin calling convention. \\
WOTS+                 & $+$262 & Defensive chain-length bound checks and the ADRS-by-pointer convention. \\
SHAKE module          & $+$458 & Incremental \emph{init / absorb / finalize / squeeze} API, instead of a single caller-allocated whole-message buffer; necessitated by the no-VLA rule. \\
\bottomrule
\end{tabularx}
\end{table}

None of these deltas reflects any property of the XMSS algorithm.
The reference's terseness comes from VLAs, contiguous
caller-allocated hash-input buffers, and 32-bit-word ADRS structures
serialised inside each hash wrapper, all incompatible with the
Jasmin subset. The headline gap of approximately 1942 instructions
between the two archives is almost exactly the size of this
project's bundled SHA-2 implementation (1953 instructions); the
upstream reference delegates SHA-256 and SHA-512 to OpenSSL's
\texttt{libcrypto} and so contains no SHA-2 internals at all. Once
SHA-2 is netted out, the two archives agree to within eleven
instructions in total size.

\section{Re-implementing XMSS in Jasmin}

The Jasmin language, by design, does not let one write the same kind
of C that the upstream reference is written in. There is no heap
allocation, no recursion, no \texttt{void *}, no function pointers,
and, as became apparent only during implementation, no way at all
to obtain a register-typed pointer to a stack-local variable.
Several patterns the C implementation takes for granted required
architectural workarounds.

\subsection{The caller-provided scratch convention}

The most consequential workaround is the \emph{caller-provided
scratch buffer}. The address structure (ADRS) is held as eight 32-bit
words on the stack and serialised to thirty-two bytes for hashing.
In C, that buffer is a stack-local. In Jasmin, the inability to take
the address of a stack-local as a register pointer forces every
function that needs ADRS bytes to receive a pointer to scratch space
from its caller. This propagates up the call graph: every
algorithm-layer function takes an explicit scratch pointer, and the
top-level export functions take a 2240-byte caller-allocated scratch
buffer.

Once adopted consistently, the convention eliminates any hidden
allocation and makes stack frames predictable, properties that
Jasmin code review and (eventually) EasyCrypt proofs benefit from.

\subsection{Compile-time parameter specialisation}

Jasmin has no const pointers and no inline parameter structs. Each
parameter set therefore compiles to its own specialised assembly
module, with parameter constants ($N$, $W$, $H$, $D$, chain length,
index width, and so on) declared in a thin top-level wrapper file
that includes the shared algorithm sources. Adding a new parameter
set is, in the present design, on the order of fifteen lines of
constant declarations and a new wrapper, with no algorithm-layer
change.

This monomorphic compilation is also a testing advantage. Each
Jasmin test exercises exactly one parameter combination through the
production code path, with no chance of a parameter being silently
substituted at runtime.

\subsection{Constant-time discipline}

The \texttt{jasmin-ct} checker verifies, at the source level, that
no value annotated as secret reaches a branch condition or a memory
address. The implementation passes the checker on all twenty-four
export functions across all five parameter sets. The checker is
strict in ways that traditional C analysers are not: it rejects, for
example, a comparison whose result is later used to index a table,
even if the table is constant.

A small number of declassification sites are needed. The principal
one is the leaf index used during signing: it is secret until
written into the signature, public thereafter, and the exhaustion
check is a branch. WOTS+ chain lengths are derived from
the message hash, which is public, and so do not raise CT
obligations for the variable-trip chain loops.

The constant-time guarantee is a source-level property of the
Jasmin program, preserved by the Jasmin compiler across backends.
It will therefore carry to RISC-V, once the Jasmin RISC-V backend
matures, without any reverification of the algorithm-layer code.

\section{RISC-V ISA analysis}

To characterise the ISA surface XMSS exercises on RISC-V, we
profiled both this project's C implementation and the upstream
reference. Both were cross-compiled with
\texttt{riscv64-linux-gnu-gcc}~13.3.0 at \texttt{-O3} for both
\texttt{rv64gc} (the standard general-purpose profile) and
\texttt{rv64gc\_zbb} (the same profile with the Zbb
bit-manipulation extension enabled). Every instruction in each
archive was classified by mnemonic against a lookup table generated
from the \texttt{riscv-opcodes} database. Object-file granularity
was preserved so that hash-layer and algorithm-layer instruction
counts could be separated.

\subsection{Required ISA: I + M is sufficient}

Both implementations, compiled for the standard \texttt{rv64gc}
profile, use only the \textbf{I} and \textbf{M} extensions.
Table~\ref{tab:isa-headline} gives the headline counts.

\begin{table}[ht]
\centering
\caption{Headline instruction counts with \texttt{-march=rv64gc -O3}.}
\label{tab:isa-headline}
\begin{tabular}{@{}lrrr@{}}
\toprule
Metric                & This project's C & Upstream reference & Difference \\
\midrule
Object files          & 13   & 9    & $-$4 \\
Total instructions    & 9562 & 7620 & $-$1942 \\
\quad I               & 9505 (99.40\%) & 7541 (98.96\%) & $-$1964 \\
\quad M               &   57 (0.60\%)  &   79 (1.04\%)  & $+$22  \\
A, F, D, Zb*          & 0    & 0    & --- \\
C-encoded ratio       & 48\% & 53\% & $+$5\,pp \\
Unique mnemonics      & 52   & 54   & $+$2 \\
\bottomrule
\end{tabular}
\end{table}

The remaining 0.6\% (1.04\% for the reference) are M-extension
multiply/divide instructions, all compiler-generated from ordinary
arithmetic on array indices and parameter values. No A, F, D, Zba,
Zbb, Zbs, Zicsr, or Zifencei instructions are emitted by either
codebase under \texttt{rv64gc}. The minimum functional ISA target
for XMSS is therefore \textbf{rv64im}. The C extension (compressed
encoding) halves average instruction width (around half of each
archive uses 16-bit encoding) but adds no semantic capability and
is invisible to source-level Jasmin code in any case.

The composition of the M-extension uses also differs between
implementations (Table~\ref{tab:m-mix}). The reference's larger
share comes mostly from inline byte-offset multiplications inside
its single combined core file and from divide/remainder operations
inside its bundled Keccak module's variable-length absorb/squeeze
loops. Neither pattern indicates structural dependence on M: an
RV64I-only target would compile both codebases unchanged with
shift-and-add expansions in place of the multiplies.

\begin{table}[ht]
\centering
\caption{M-extension mnemonic breakdown.}
\label{tab:m-mix}
\begin{tabular}{@{}lrrl@{}}
\toprule
Mnemonic        & This project & Reference & Compiler emits it for \\
\midrule
\texttt{mulw}   & 37 & 68 & 32-bit index $\times$ parameter \\
\texttt{mul}    & 17 &  6 & 64-bit offset / pointer arithmetic \\
\texttt{divuw}  &  3 &  1 & parameter derivation \\
\texttt{divu}   &  0 &  2 & SHAKE absorb/squeeze offset (reference only) \\
\texttt{remu}   &  0 &  2 & SHAKE block-length modulus (reference only) \\
\midrule
Total           & 57 & 79 & \\
\bottomrule
\end{tabular}
\end{table}

\subsection{Per-module breakdown}

The two codebases carve XMSS into modules differently. The reference
bundles WOTS+ public-key generation, L-tree, treehash, BDS state
management, and XMSS / XMSS-MT signing into a single core file; this
project splits those concerns across separate compilation units.
Table~\ref{tab:per-module} shows representative per-module counts.
The reference's hash module is conspicuously small (684 instructions
with no M uses) because it is a thin shim that delegates SHA-256 and
SHA-512 to OpenSSL's \texttt{libcrypto}; SHA-2 internals are not
present in the analysed reference archive at all.

\begin{table}[ht]
\centering
\caption{Per-module instruction counts under \texttt{-march=rv64gc}.
  Modules are grouped by role; the reference's combined core file is
  shown alongside this project's split modules.}
\label{tab:per-module}
\begin{tabular}{@{}lrr@{}}
\toprule
Role / module                              & This project & Reference \\
\midrule
\multicolumn{3}{@{}l}{\textit{Hash layer}} \\
\quad SHA-2 core                           & 1953  & ---\,(in \texttt{libcrypto}) \\
\quad Keccak / SHAKE                       & 1337  & 1259 \\
\quad Hash wrappers                        & 1082  &  684 \\
\midrule
\multicolumn{3}{@{}l}{\textit{Algorithm layer}} \\
\quad WOTS+                                &  760  &  498 \\
\quad L-tree + treehash + BDS runtime      & 1743  & $\subset$\,3169 \\
\quad BDS serialisation                    &  473  & ---\,(not present) \\
\quad XMSS / XMSS-MT core                  & 1782  & 3682 \\
\quad ADRS manipulation                    &   99  &   29 \\
\quad Parameters                           &  286  & 1077 \\
\quad Utilities                            &   47  &   29 \\
\midrule
Total                                      & 9562  & 7620 \\
\bottomrule
\end{tabular}
\end{table}

The headline 1942-instruction archive-size gap is almost entirely
the size of this project's bundled SHA-2 implementation
(1953 instructions). Once SHA-2 is netted out, the two archives
agree to within eleven instructions in total size. The remaining
per-module deltas---hash wrappers, ADRS, WOTS+, parameters---are
the design-tax items already discussed in
Table~\ref{tab:tax}.

\subsection{Where Zbb lands}

Re-compiling with \texttt{-march=rv64gc\_zbb} introduces 164 Zbb
instructions in this project's archive (132 in the reference's), as
shown in Table~\ref{tab:zbb-headline}. Both archives shed base-integer
instructions (Zbb collapses multi-instruction rotation, $\overline{a}
\land b$, and branch-based min/max sequences into single ops); the
reference saves more in absolute terms because its baseline is
smaller and its compositional structure exposes more
recognisable idioms (see \S\ref{sec:zbb-asymmetry}).

\begin{table}[ht]
\centering
\caption{Effect of enabling Zbb (\texttt{-march=rv64gc\_zbb}).}
\label{tab:zbb-headline}
\begin{tabular}{@{}lrrrrrr@{}}
\toprule
                & \multicolumn{2}{c}{This project's C} & \multicolumn{2}{c}{Reference} \\
\cmidrule(lr){2-3}\cmidrule(lr){4-5}
                & rv64gc & rv64gc\_zbb & rv64gc & rv64gc\_zbb \\
\midrule
I               & 9505 & 9281  & 7541 & 6982 \\
M               &   57 &   57  &   79 &   79 \\
Zbb             &    0 &  164  &    0 &  132 \\
Total           & 9562 & 9502  & 7620 & 7193 \\
\midrule
Net change      &      & $-$60 &      & $-$427 \\
\bottomrule
\end{tabular}
\end{table}

The mnemonic mix (Table~\ref{tab:zbb-mix}) is structurally
asymmetric. This project's archive sees substantial 32-bit SHA-2
rotations (\texttt{rolw}, \texttt{roriw}) and byte-reversals
(\texttt{rev8}); the reference's archive sees almost none, because
its SHA-2 implementation is not in the archive (delegated to
\texttt{libcrypto}). Conversely, the reference shows more 64-bit
Keccak rotations (\texttt{rol}, \texttt{rori}) and significantly more
\texttt{andn} uses, reflecting the prominence of the Keccak $\chi$
step in its bundled SHAKE module.

\begin{table}[ht]
\centering
\caption{Zbb mnemonic mix with \texttt{-march=rv64gc\_zbb}.}
\label{tab:zbb-mix}
\begin{tabular}{@{}lrrlp{3.5cm}@{}}
\toprule
Mnemonic        & This project & Reference & Replaces & Where \\
\midrule
\texttt{rolw}   & 41 &  0 & 3 insns & 32-bit SHA-2 rotation \\
\texttt{roriw}  & 17 &  0 & 3 insns & 32-bit SHA-2 rotation \\
\texttt{rev8}   & 19 &  0 & multi-insn byte-swap & SHA-2 endianness \\
\texttt{rol}    & 17 & 38 & 3 insns & 64-bit Keccak rho \\
\texttt{rori}   & 27 & 20 & 3 insns & 64-bit Keccak rho \\
\texttt{andn}   & 28 & 51 & 2 insns & Keccak chi, masks \\
\texttt{maxu}   & 12 & 20 & branch-based & BDS height comparisons \\
\texttt{minu}   &  3 &  2 & branch-based & BDS / SHAKE min \\
\texttt{zext.h} &  0 &  1 & 2 insns & 16-bit widening \\
\midrule
Total           & 164 & 132 & & \\
\bottomrule
\end{tabular}
\end{table}

Both implementations show the same qualitative pattern:
\textbf{Zbb uptake concentrates almost entirely in the hash layer.} In this project's archive 142 of 164 Zbb
instructions (87\%) land inside SHA-2 and SHAKE; in the reference,
108 of 132 (82\%) land inside its bundled Keccak. The handful of Zbb
uses that escape to the algorithm layer (around 22--24 in each
implementation) are opportunistic \texttt{maxu}/\texttt{minu}/\texttt{andn}
applications for index comparisons and bitmask combinations. None
of these is structural: the algorithm runs equivalently without
them.

\subsection{Why the absolute Zbb savings are asymmetric}
\label{sec:zbb-asymmetry}

The two implementations \emph{benefit} from Zbb to very different
degrees: this project's archive nets 60 instructions saved; the
reference nets 427. Two effects combine. First, this
project's archive contains 32-bit SHA-2 rotations and byte-reversals
that the reference does not, but those instructions replace
three-instruction software sequences with a single instruction, so
the proportional saving is bounded. Second, this project's SHAKE
module is larger than the reference's by approximately 78
instructions because of its incremental context-based API
(\emph{init / absorb / finalize / squeeze}, required by the no-VLA
rule), whereas the reference uses a single contiguous-buffer call
shape; this raises the baseline against which the absolute saving
is measured. A reference build that included SHA-2 in source form
would show a substantially larger absolute Zbb benefit, dominated
by the SHA-2 compression core's rotations and byte-reversals.

\subsection{Implication for ISA-tuned code generation}

The implication for any future ISA-tuning work---whether on the
Jasmin RISC-V backend, on hand-written assembly intrinsics, or on
the choice of pre-built primitive libraries---is that the payoff
sits in the hash layer. The algorithm layer (WOTS+, L-tree,
BDS, treehash, XMSS, XMSS-MT) is load--store--branch--shift code
with no structural dependence on any extension beyond I~+~M, and
ISA-specific work there will not move performance.

\section{Scoping a future Jasmin RISC-V port}
\label{sec:scoping}

The ISA findings above bear on a port-scoping question: when a
Jasmin RISC-V port becomes feasible, should it cover the whole
stack, or only the hash primitives?

The default assumption is that the existing Jasmin algorithm-layer
code should be recompiled for RISC-V. This path has appeal: the
constant-time properties already established on x86-64 are
source-level Jasmin properties and would carry to RISC-V without
re-verification, and no new algorithm-layer work would be required.

The ISA evidence complicates the picture. Three observations bear
on the decision:

\begin{enumerate}[nosep]
  \item Zbb uptake concentrates in the hash layer ($\geq 82\%$ in
    both codebases). The algorithm layer derives no ISA-specific
    benefit from a Jasmin port.
  \item Upper-layer code-size differences between this project's C
    and the upstream reference are entirely Jasmin-portability tax
    (Table~\ref{tab:tax}). A Jasmin algorithm layer pays this tax in
    source size and compilation overhead without buying any
    RV64-specific advantage.
  \item Algorithm-layer code is ISA-neutral. A C or Rust upper
    layer over a Jasmin hash primitive would compile to substantially
    similar instruction streams.
\end{enumerate}

The split-layer alternative this evidence suggests is to keep
Jasmin for the hash primitives (where it earns formally checked
constant-time \emph{and} ISA-specific code generation) but write
the algorithm layer in a less-restricted language: C, or, more
attractively, Rust, which adds memory-safety guarantees that the
algorithm layer cannot otherwise express. The precise integration
shape (header-based FFI, language bindings, build-time composition)
is an open question and is not prescribed here.

The alternative gives up three things:

\begin{itemize}[nosep]
  \item \textbf{End-to-end constant-time guarantees through
    Jasmin.} The \texttt{jasmin-ct} property is a source-level
    property of Jasmin code; once control crosses an FFI boundary
    into a different language, any CT property of the calling code
    has to be argued separately. The loss is bounded: the algorithm
    layer's CT obligations
    (WOTS+ chain index, BDS array indexing) are simple enough to
    audit by inspection or with tools targeting the host language,
    and the hard CT cases (variable-time instructions inside hash
    rounds) all live in the hash layer, which remains under Jasmin.
  \item \textbf{The finished Jasmin upper layer.} The existing
    Jasmin algorithm-layer code is complete, tested, CT-verified,
    and cross-validated against the C reference. Replacing it would
    discard a finished body of work; this is a real cost even if
    the forward case for the split-layer architecture is strong.
  \item \textbf{Symmetry with a future formal-verification effort.}
    If the project later pursues a full EasyCrypt-style correctness
    proof, that proof infrastructure is for Jasmin code; an upper
    layer in another language puts an EasyCrypt proof of the upper
    layer out of reach. No such proof effort is currently underway.
\end{itemize}

This section records the evidence rather than recommending a path;
the decision needs both the ISA argument and the integrity
arguments side by side.

\section{Known limitations and open work}

Both implementations are functionally complete within their
declared scope, but several known limitations remain. At the
time of writing, \textbf{no security-affecting defects remain
outstanding.} Four
security items previously identified during review---secret-material
zeroing across hash and signing paths; reliance on volatile
semantics in constant-time comparison and memory-zeroing utilities;
missing bounds validation on the BDS retain parameter; and toolchain
permission hardening---have all been resolved during the work
covered here. The limitations that remain fall into three
categories.

\subsection{Functional and specification issues}

\begin{itemize}[nosep]
  \item \textbf{Jasmin verify does not validate the OID embedded in
    the public key}, an explicit deviation from the Jasmin
    implementation's own specification document, which states that
    verify must check it. The omission is partially mitigated by
    Jasmin's compile-time parameter specialisation: a public key
    with a mismatched OID would already fail verification at the
    root-comparison step, so the security impact is a confusing
    error message rather than an unsound accept. The C
    implementation does perform the check.
  \item \textbf{XMSS and XMSS-MT OID namespaces overlap.}
    RFC~8391 assigns separate OID registries for XMSS and XMSS-MT;
    the value \texttt{0x00000001} is XMSS-SHA2\_10\_256 in one
    registry and XMSSMT-SHA2\_20/2\_256 in the other. A caller who
    passes a public key into the wrong verify entry point with a
    parameter struct that happens to share the same OID value will
    get a spurious OID-check pass and a downstream signature-shape
    failure. This is an RFC design limitation rather than an
    implementation bug; the OID check on its own does not
    distinguish XMSS from XMSS-MT keys.
\end{itemize}

\subsection{Coverage gaps in the test suite}

\begin{itemize}[nosep]
  \item Cross-implementation deep interop is currently exercised on
    the multi-tree 20/2 parameter set but not on 20/4. The 20/4
    set is exercised by the C exhaustive boundary-crossing test
    and by the Jasmin API roundtrip, but there is no per-signature
    BDS-state comparison between implementations for it.
  \item The Jasmin implementation lacks an explicit key-exhaustion
    test (signing through the entire $2^H$ index space and
    confirming the documented refusal at exhaustion). The behaviour
    is exercised indirectly by the C implementation's exhaustive
    tests but should be verified directly on the Jasmin side.
  \item The C unit-test suite does not cover the L-tree and treehash
    modules in isolation. Both are exercised by every higher-level
    test, but a targeted unit test would isolate regressions.
  \item The C test suite's coverage of the SHAKE backend and of the
    $n=64$ parameter sets is thinner than its SHA-2/256 coverage.
  \item The XMSS-MT unit test only spot-checks four of the 1025
    signatures around a tree boundary; full per-signature coverage
    of the boundary would catch a wider class of state-machine
    regressions.
\end{itemize}

\subsection{Code-quality and ergonomic issues}

A number of cleanup and ergonomic items are tracked: a wasteful dual
hash-context declaration in one C hash-message wrapper; absence of
a signature-length parameter on verify, which would allow callers to
detect truncated inputs; subtle BDS swap logic in XMSS-MT
sign that would benefit from documentation; a moderate stack-pressure
concern (a 34~KB stack-allocated swap temporary in XMSS-MT sign,
which is fine on hosted targets but undesirable for embedded use);
and helper-function duplication across Jasmin test files.

\subsection{Scope work that has not been started}

Three substantive pieces of work are explicitly out of scope for
the implementations as they stand:

\begin{itemize}[nosep]
  \item \textbf{Hash backends beyond SHA-2/256 in Jasmin.} SHA-2/512
    and SHAKE-128/256 backends are designed but not written;
    completing them would extend Jasmin coverage to the remaining
    parameter sets registered in RFC~8391.
  \item \textbf{Code review of the Jasmin sources.} A full review
    pass against the Jasmin specification document is open.
  \item \textbf{Formal verification beyond constant-time.} The
    EasyCrypt proof effort for functional correctness has not begun.
\end{itemize}

\section{Methodological observations}

A few observations from the implementation process transfer to the
next verification-oriented re-implementation effort and are not
peculiar to XMSS.

\paragraph{Specifications are most useful when they are written
independently of the code.}
The Jasmin implementation's specification document was written from
the RFC and from declared design decisions, deliberately without
reference to the Jasmin source. This makes it a check on the
implementation rather than a description of it. The verify-side
OID-check deviation described in the previous section was identified
by a routine spec-versus-code reading.

\paragraph{Stateful-scheme test discipline must be exhaustive at
small parameter sizes.}
The natural failure mode of a stateful authentication-path scheme
is a state-machine bug that manifests only at specific index
boundaries and only on traversals that cross them. Spot-check tests
that sign one or two messages and verify them are insufficient to
distinguish ``the signing function works'' from ``the BDS state
update preserves the invariants required for the next signature.''
Exhaustive sign/verify across every leaf index for valid
small-height parameter combinations is tractable and has caught
concrete bugs in this project that wider but shallower test
coverage missed. Cross-implementation interop should compare
internal BDS state byte-for-byte after every signature, not just
the signature outputs.

\paragraph{Test each layer in isolation before building on it.}
The Jasmin implementation grew bottom-up: address structure,
utilities, SHA-256, WOTS+, L-tree, treehash, BDS, XMSS, XMSS-MT,
each tested against the C output before the next layer was written.
The top-level XMSS export functions assembled with no
algorithmic bugs, because every component had already been
validated. By contrast, attempts to write multiple top-level
functions in one pass before compiling produced tangled
register-pressure failures that took considerably longer to unbraid.

\paragraph{Reach for runtime tools first when the failure is a
runtime symptom.}
A non-determinism in early Jasmin keygen output was diagnosed in
seconds by a memory-error detector. The same bug was opaque to
manual inspection of the surrounding register-spill-heavy assembly.
Code review is the right tool for compile-time errors and design
issues; for wrong output, non-determinism, or crashes, the runtime
tool should come first.

\paragraph{A previous claim that ``X passes'' is unverified until
the test that demonstrates it is named.}
Several incidents in this project trace back to summary claims of
correctness that, on inspection, were not backed by tests covering
the parameter combination in question. The corollary is the test
discipline now in place: every parameter set is signed through at
least one full tree boundary, and any claim of correctness names the
test that demonstrates it.

\section{Summary of conclusions}

\begin{description}[style=nextline,leftmargin=1.5cm]
  \item[Two complete, cross-validated implementations.]
    A C99 reference covering all RFC~8391 parameter sets and a
    Jasmin implementation covering five SHA-2/256 parameter sets
    pass exhaustive cross-validation against the upstream reference
    and against each other. The Jasmin implementation is verified
    constant-time on x86-64 by \texttt{jasmin-ct}.

  \item[XMSS uses only RV64I + M.]
    Both implementations compile to base integer plus a small
    handful of compiler-emitted multiply instructions on
    \texttt{rv64gc}. The minimum functional ISA target is
    \texttt{rv64im}. No floating-point, atomic, or bit-manipulation
    extensions are required.

  \item[Zbb lives in the hash layer.]
    Where the Zbb extension is enabled, $\geq 82\%$ of the
    additional Zbb instructions land in SHA-2 / SHAKE / Keccak.
    The algorithm layer benefits negligibly. Any ISA-tuning effort
    for XMSS---compiler, hand-written assembly, or backend
    development---should target the hash primitives.

  \item[The algorithm layer is design-tax-bound, not algorithm-bound.]
    Module-level instruction count differences between this
    project's C and the upstream reference are explained entirely by
    Jasmin-portability constraints. The XMSS algorithm itself is
    small.

  \item[The upper-layer scoping question is open.]
    The existing Jasmin upper layer is finished and tested; the ISA
    evidence suggests its constraints buy little RV64-specific
    benefit. A split-layer alternative (Jasmin hash primitives,
    Rust or C upper layer) is plausible; the trade-offs are
    discussed in \S\ref{sec:scoping} without prescribing a choice.

  \item[No security-affecting defects currently outstanding.]
    Four security items identified during review have all been
    resolved. Remaining limitations are concentrated in
    testing-coverage gaps, code-quality items, and unfinished scope
    (hash backends beyond SHA-2/256 in Jasmin). One functional
    deviation is documented: Jasmin verify does not validate the
    public-key OID, though compile-time parameter specialisation
    means a mismatched OID would still fail verification at the
    root-comparison step.
\end{description}

\begin{thebibliography}{9}
  \bibitem{BDS09}
  Johannes Buchmann, Erik Dahmen, and Andreas H\"{u}lsing.
  \emph{XMSS---A practical forward secure signature scheme based on
    minimal security assumptions.}
  In Post-Quantum Cryptography (PQCrypto 2011), LNCS 7071, 2011.
\end{thebibliography}

\end{document}
